% ==============================================================================
% PAGE 34: CHAPTER-4 FABRICATION - PART 1 (Numbered Page 26)
% ==============================================================================
\begin{center}
  {\fontsize{15}{18}\selectfont \bfseries \textcolor{red!80!black}{\underline{CHAPTER-4}}}\\[2mm]
  {\fontsize{16}{20}\selectfont \bfseries \textcolor{red!80!black}{\underline{FABRICATION}}}\\[5mm]
\end{center}

\subsubsection*{1. Frame Construction}
\noindent
\begin{spacing}{1.2}
{\fontsize{10.2}{14.5}\selectfont
The base frame is made of wooden sheet ($60 \times 60$~cm) to provide stability and support to the setup. Holes are drilled to fix the shaft supports at the center.\\[2.5mm]
The frame is painted by white colour because it reflect solar rays to prevent electronic component from over heating.
}
\end{spacing}

\vspace{4mm}
\subsubsection*{2. Parabolic Dish Fabrication}
\noindent
\begin{spacing}{1.2}
{\fontsize{10.2}{14.5}\selectfont
A parabolic dish is made using satellite dish for lightweight and good reflection.\\[2.5mm]
The dish surface is covered with reflective mirror pieces to concentrate sunlight at the focus point. The dish is mounted on a PVC pipe shaft that allows rotational motion.
}
\end{spacing}

\vspace{4mm}
\subsubsection*{3. Receiver Fabrication}
\noindent
\begin{spacing}{1.2}
{\fontsize{10.2}{14.5}\selectfont
The receiver pipe (usually copper tube) is placed at the focal point of the parabola dish.\\[2.5mm]
The copper tube is black-painted for maximum heat absorption. The tube is connected with inlet and outlet pipes to circulate water.\\[2.5mm]
Copper tube is placed on the insulating material (aluminium plate and glass wool) for reducing heat loss.
}
\end{spacing}

\vspace{4mm}
\subsubsection*{4. Tracking Mechanism}
\noindent
\begin{spacing}{1.2}
{\fontsize{10.2}{14.5}\selectfont
A single-axis solar tracking system is used to follow the sun's motion east to west.\\[2.5mm]
A PMDC motor with pulley and V-belt attached with the PVC pipe shaft which is used to rotate the parabolic dish. Sensors (LDRs) are placed on both sides of the dish to detect sunlight intensity difference.\\[2.5mm]
A motor driver circuit and microcontroller (Arduino nano) control the motor rotation based on sensor signals.
}
\end{spacing}
\clearpage
